\exnumber{00}
\extitle{UI V - Ex 1 - Audit}
\exfiles{ex00/archive zip, nomme Name\_FirstName\_DeliverableName\_Date\_VersionNumber}
\exauthorize{Ceux nommés dans l'exercice}
\makeheaderfiles

\textbf{Contexte :}

Bienvenue dans cette toute dernière quête !

L'objectif est d'endosser le rôle d'auditeur UX. On vous demandera de réaliser l'audit d'un site web ou d'une application, puis de rédiger une série de recommandations pour l'améliorer.

Soyons critiques, créatifs et impactants !

\textbf{Instructions :}

Réalisez un audit heuristique sur l'un de ces 3 sites web ou applications différents :

\begin{itemize}
\item https://www.ikea.com/ \url{https://www.ikea.com/}
\item https://www.doctorswithoutborders.org/ \url{https://www.doctorswithoutborders.org/}
\item Application WikiArt
\end{itemize}

\textbf{Livrables :}

\begin{itemize}
\item Un graphique complété
\item Un rapport comprenant les 5 étapes recommandées décrites ci-dessous
\item Déposez les documents sur Github, dans un zip nommé « ProjectTitle\_Name\_FirstName », et chaque livrable ainsi : « Name\_FirstName\_DeliverableName\_Date\_VersionNumber ». Exemple : « Doe\_John\_MidFiPrototype\_05242024\_V1 ».
\item Ajoutez un titre dans tous les documents écrits.
\item Soignez l'aspect global des livrables. Ils doivent être clairs, simples et lisibles. Vous pouvez regarder canva.com, sans surcharger le dessin de détails.
\end{itemize}

\textbf{Pistes :}

\begin{itemize}
\item Étapes recommandées
\begin{itemize}
\item Définir le périmètre.
\end{itemize}
\end{itemize}
1. Connaître les exigences business et le profil démographique des utilisateurs finaux.
2. Décidez quels outils de reporting et quelles heuristiques utiliser (dans ce cas : cette checklist d'évaluation heuristique \url{https://drive.google.com/file/d/10KbfbNZA1oVS1sXbjjXLPPmdZ6nqVkdc/view}).
3. Évaluer l'expérience et identifier les problèmes d'utilisabilité.
4. Analyser et agréger les données.
5. Présenter les résultats.

\textbf{Ressources :}

\begin{itemize}
\item Introduction à l'évaluation heuristique : un guide pour débutants \url{https://uxdesign.cc/introduction-to-heuristic-evaluation-658705606518}
\item Les 10 heuristiques d'utilisabilité de Nielsen \url{https://www.nngroup.com/articles/ten-usability-heuristics/}
\item Ce que vous retirez vraiment d'une évaluation heuristique \url{https://uxmag.com/articles/what-you-really-get-from-a-heuristic-evaluation}
\item Analyse heuristique pour l'UX : comment mener une évaluation d'utilisabilité \url{https://uxdesign.cc/heuristic-analysis-for-ux-how-to-run-a-usability-evaluation-12c86d43936f}
\item L'analyse heuristique dans le processus de design \url{https://uxdesign.cc/heuristic-analysis-in-the-design-process-usability-inspection-methods-d200768eb38d}
\item Une nouvelle checklist d'évaluation heuristique d'utilisabilité \url{https://uxplanet.org/a-new-usability-heuristic-evaluation-checklist-259f588da308}
\item Outil : checklist d'évaluation heuristique \url{https://drive.google.com/file/d/10KbfbNZA1oVS1sXbjjXLPPmdZ6nqVkdc/view}
\end{itemize}

